% !TEX root = ./rapport.tex
% ══════════════════════════════════════════════════════════════════════
%  CHAPITRE 6 — SÉCURITÉ DE LA PLATEFORME SYNAPSE
% ══════════════════════════════════════════════════════════════════════

\chapter{Sécurité de la Plateforme}

\noindent\textbf{Plan}
\begin{enumerate}
  \item Introduction
  \item Architecture de sécurité globale
  \item Authentification OAuth2/OpenID Connect avec Keycloak
  \item Contrôle d’accès basé sur les rôles (RBAC)
  \item Sécurisation des microservices Spring Boot
  \item Sécurisation du canal WebSocket STOMP
  \item Prévention des vulnérabilités OWASP
  \item Gestion des données sensibles
  \item Traçabilité et journalisation des audits
  \item Conclusion
\end{enumerate}

% ─────────────────────────────────────────────────────────────
\section{Introduction}
% ─────────────────────────────────────────────────────────────

Une plateforme de gestion des ressources humaines concentre par nature des données
hautement sensibles~: identités, salaires, données médicales (congés longue maladie,
comité médical), informations contractuelles et dossiers personnels. La sécurité de
Synapse n’est pas une couche ajoutée \textit{a posteriori}~: elle constitue une
contrainte transversale intégrée à chaque décision architecturale, depuis la
conception du schéma Oracle jusqu'à la validation des tokens JWT dans les
intercepteurs STOMP.

Ce chapitre décrit l’ensemble des mécanismes de sécurité implémentés dans Synapse,
de l’authentification centralisée Keycloak à la prévention des vulnérabilités OWASP,
en passant par le contrôle d’accès à granularité fine et la sécurisation des
communications temps réel.

% ─────────────────────────────────────────────────────────────
\section{Architecture de sécurité globale}
% ─────────────────────────────────────────────────────────────

\subsection{Périmètre et surface d’attaque}

La plateforme Synapse expose les surfaces suivantes, chacune associée à des contrôles
de sécurité spécifiques~:

\begin{table}[H]
\centering
\caption{Surfaces d'exposition et mécanismes de sécurité associés}
\label{tab:surfaces_securite}
\begin{tabularx}{\textwidth}{VlVLVLV}
\rowcolor{headerblue!55}
\hline
\textbf{Surface} & \textbf{Description} & \textbf{Contrôle de sécurité} \\
\hline
API REST (HTTP) & Endpoints exposés par les 7 microservices & JWT Keycloak + \texttt{@PreAuthorize} \\
\hline
WebSocket (STOMP) & Canal temps réel du chat-service & Validation JWT dans le STOMP interceptor \\
\hline
Frontend SPA & Application Angular 21 & Guard de routage + PKCE OAuth2 \\
\hline
Base de données & Oracle 23ai (schéma GERAI) & Accès via HikariCP (compte dédié) \\
\hline
Keycloak Admin API & Provisioning des utilisateurs & Admin Client credentials (backend uniquement) \\
\hline
API Groq (externe) & Appels LLM (chatbot, scoring CV) & Clé API côté serveur uniquement \\
\hline
\end{tabularx}
\end{table}

\subsection{Modèle Zero Trust appliqué}

Synapse adopte un modèle \textit{Zero Trust}~: aucun microservice n’est considéré
comme fiable par défaut, même s’il appartient au réseau interne Docker. Les appels
inter-services REST réutilisent le JWT de l’utilisateur connecté, transmis en
en-tête \texttt{Authorization~: Bearer} lors de chaque requête. L’isolation réseau
est assurée par Docker Compose~: les conteneurs des microservices ne sont pas
exposés directement sur l’interface publique~; seuls les ports applicatifs nécessaires
sont mappés sur \texttt{localhost}.

Les contrôles de sécurité s’articulent en trois périmètres concentriques~:
(1)~le \textbf{périmètre réseau}, où Docker Compose isole les conteneurs sur un
réseau interne privé~; (2)~le \textbf{périmètre applicatif}, où chaque requête
HTTP et WebSocket est validée par un JWT Keycloak avant traitement~; (3)~le
\textbf{périmètre données}, où Oracle applique les contraintes d’intégrité et
HikariCP filtre les connexions via un compte dédié à droits limités.

% ─────────────────────────────────────────────────────────────
\section{Authentification OAuth2/OpenID Connect avec Keycloak}
% ─────────────────────────────────────────────────────────────

\subsection{Flux d’authentification PKCE}

Synapse implémente le flux \textbf{Authorization Code avec PKCE}
(\textit{Proof Key for Code Exchange}), recommandé par l'IETF (RFC 7636) pour les
applications SPA car il ne stocke jamais de \textit{client secret} dans le code
JavaScript. Le flux se déroule en dix étapes documentées dans le diagramme de
séquence du chapitre 3 (figure~\ref{fig:seq_auth}).

\textbf{Paramètres de sécurité du flow PKCE~:}
\begin{itemize}[leftmargin=*,nosep]
  \item \texttt{code\_challenge\_method}~: \texttt{S256} (SHA-256)
  \item \texttt{access\_token\_lifespan}~: 5 minutes (par défaut Keycloak)
  \item \texttt{refresh\_token\_lifespan}~: 30 minutes
  \item \texttt{pkce\_enabled}~: \texttt{true} dans la configuration du client Keycloak
\end{itemize}

\subsection{Structure du JWT émis par Keycloak}

Chaque \textit{access token} JWT émis par Keycloak contient les claims suivants,
exploités par les microservices Spring Boot~:

\begin{table}[H]
\centering
\caption{Claims JWT utilisés dans Synapse}
\label{tab:jwt_claims}
\begin{tabularx}{\textwidth}{VlVlVLV}
\rowcolor{headerblue!55}
\hline
\textbf{Claim} & \textbf{Type} & \textbf{Utilisation dans Synapse} \\
\hline
\texttt{sub} & UUID & Identifiant Keycloak~; résolu en \texttt{employee\_id} Oracle via \texttt{findByKeycloakSub()} \\
\hline
\texttt{realm\_access.roles} & Tableau & Rôles organisationnels~; extraits par \texttt{KeycloakJwtRoleConverter} \\
\hline
\texttt{preferred\_username} & String & Nom d'utilisateur affiché dans l’interface \\
\hline
\texttt{email} & String & E-mail de l’utilisateur (synchronisé avec Oracle) \\
\hline
\texttt{exp} & Timestamp & Date d’expiration~; vérifiée automatiquement par Spring Security \\
\hline
\texttt{iss} & URI & Issuer URI Keycloak~; validé contre \texttt{issuer-uri} configuré \\
\hline
\end{tabularx}
\end{table}

\subsection{Configuration keycloak-angular côté frontend}

L’intégration Keycloak dans Angular est réalisée via \texttt{keycloak-angular 21}.
\textit{Note~: les identifiants techniques (\texttt{realm~: gerai}, \texttt{clientId~: gerai-app},
schéma Oracle \texttt{GERAI}) conservent le nom initial du projet~; la dénomination
commerciale «~Synapse~» a été adoptée en cours de développement sans renommage des artefacts techniques.}
La bibliothèque initialise le client OIDC avant le démarrage de l’application,
injecte automatiquement le Bearer token dans chaque requête HTTP via
\texttt{HttpClientModule}, et gère le renouvellement silencieux du token via
\texttt{refresh\_token}.

\begin{table}[H]
\centering
\caption{Configuration keycloak-angular dans Synapse}
\label{tab:kc_angular_config}
\begin{tabularx}{\textwidth}{VlVlVLV}
\rowcolor{headerblue!55}
\hline
\textbf{Paramètre} & \textbf{Valeur} & \textbf{Rôle} \\
\hline
\texttt{realm} & \texttt{gerai} & Realm Keycloak dédié à Synapse \\
\hline
\texttt{clientId} & \texttt{gerai-app} & Client OIDC public (PKCE) \\
\hline
\texttt{onLoad} & \texttt{login-required} & Force l’authentification avant tout accès \\
\hline
\texttt{checkLoginIframe} & \texttt{false} & Désactivé~: la vérification de session par iframe est incompatible avec la politique \texttt{SameSite=Strict} des navigateurs modernes (Chrome, Firefox) \\
\hline
\texttt{pkceMethod} & \texttt{S256} & PKCE activé avec SHA-256 \\
\hline
\end{tabularx}
\end{table}

% ─────────────────────────────────────────────────────────────
\section{Contrôle d’accès basé sur les rôles (RBAC)}
% ─────────────────────────────────────────────────────────────

\subsection{Matrice des rôles et permissions}

Synapse définit six rôles Keycloak correspondant aux niveaux hiérarchiques d’Arabsoft~:
\texttt{EMPLOYE}, \texttt{CHEF}, \texttt{RH}, \texttt{ADMIN\_RH}, \texttt{DG} et
\texttt{COMITE}. Chaque rôle délimite un périmètre fonctionnel précis~: l’employé
dépose et suit ses propres demandes~; le chef valide celles de son équipe et pilote
les projets~; le service RH instruit les dossiers~; l’administrateur RH dispose des
droits les plus étendus (CRUD employés, rapports, attrition)~; le directeur général
intervient uniquement sur les décisions de prêt~; le comité de crédit se limite au
vote. Cette ségrégation garantit qu’aucun utilisateur n’accède à des données ou
des actions hors de son périmètre organisationnel. La matrice complète des droits
par fonctionnalité est présentée en Annexe~V (voir tableau~\ref{tab:rbac_matrix}).

\subsection{Implémentation du RBAC côté backend}

Le RBAC backend repose sur deux mécanismes complémentaires~:

\paragraph{1. KeycloakJwtRoleConverter.} Ce composant Spring implémente
\texttt{Converter<Jwt, Collection<GrantedAuthority>>} et extrait le tableau
\texttt{realm\_access.roles} du token JWT pour le transformer en une collection
de \texttt{GrantedAuthority} préfixés \texttt{ROLE\_}. Ce préfixe est requis par
\texttt{@PreAuthorize("hasRole('X')")} dans Spring Security.

\paragraph{2. SecurityConfig par service.} Chaque microservice définit son propre
\texttt{SecurityFilterChain} avec des règles \texttt{requestMatchers} granulaires.
Extrait de la configuration du \textit{employe-service}~:

\begin{table}[H]
\centering
\caption{Règles de contrôle d’accès du employe-service}
\label{tab:employe_acl}
\begin{tabularx}{\textwidth}{VlVlVLV}
\rowcolor{headerblue!55}
\hline
\textbf{Méthode HTTP} & \textbf{Endpoint} & \textbf{Rôles autorisés} \\
\hline
GET & \texttt{/employees/**} & \texttt{ADMIN\_RH, RH, CHEF, EMPLOYE} \\
\hline
POST & \texttt{/employees/**} & \texttt{ADMIN\_RH, RH} \\
\hline
PUT & \texttt{/employees/**} & \texttt{ADMIN\_RH, RH} \\
\hline
DELETE & \texttt{/employees/**} & \texttt{ADMIN\_RH} \\
\hline
ANY & \texttt{/admin/keycloak/**} & \texttt{ADMIN\_RH} \\
\hline
\end{tabularx}
\end{table}

\subsection{Implémentation du RBAC côté frontend}

Côté Angular, le RBAC est implémenté à deux niveaux~:

\begin{enumerate}[leftmargin=*,nosep]
  \item \textbf{Guard de routage (\texttt{requireRoleGuard})}~: chaque route lazy-loaded
        déclare dans son champ \texttt{data.roles} le tableau des rôles autorisés.
        Le guard lit les rôles du token JWT Keycloak et redirige vers \texttt{/access-denied}
        si le rôle de l’utilisateur n’est pas présent dans la liste.
  \item \textbf{Guard de redirection (\texttt{roleRedirectGuard})}~: à l'arrivée
        sur la racine \texttt{/}, ce guard lit le rôle principal du JWT et redirige
        automatiquement selon la hiérarchie de rôles~: \texttt{comite} →
        \texttt{/admin/comite-credit}~; \texttt{admin} / \texttt{admin\_rh} /
        \texttt{directeur\_general} → \texttt{/admin/dashboard}~;
        \texttt{chef} → \texttt{/chef/dashboard}~;
        \texttt{employe} → \texttt{/employe/dashboard}.
        Tout rôle non reconnu est redirigé vers \texttt{/access-denied}.
\end{enumerate}

% ─────────────────────────────────────────────────────────────
\section{Sécurisation des microservices Spring Boot}
% ─────────────────────────────────────────────────────────────

\subsection{Session stateless (JWT)}

Tous les microservices sont configurés en mode \textbf{STATELESS}~:

\begin{itemize}[leftmargin=*,nosep]
  \item \texttt{SessionCreationPolicy.STATELESS}~: aucune session HTTP serveur
        n’est créée ni utilisée.
  \item Chaque requête est auto-suffisante~: le JWT porte toutes les informations
        d’identité et de rôle nécessaires à l’autorisation.
  \item Pas de cookie de session~: élimine les attaques CSRF basées sur les
        sessions (voir section~\ref{sec:csrf}).
\end{itemize}

\subsection{Validation du JWT}

La validation du JWT est déléguée à Spring Security OAuth2 Resource Server via
la propriété \texttt{spring.security.oauth2.resourceserver.jwt.jwk-set-uri}.
Spring Security télécharge et met en cache la clé publique Keycloak depuis
l'endpoint JWKS (\textit{JSON Web Key Set}), puis vérifie~:

\begin{enumerate}[leftmargin=*,nosep]
  \item La signature RSA du token (via RS256).
  \item La valeur du claim \texttt{iss} (issuer) contre \texttt{http://localhost:8080/realms/gerai}.
  \item La date d’expiration (\texttt{exp}).
  \item L'audience (\texttt{aud}) si configurée.
\end{enumerate}

\subsection{Configuration CORS}

Chaque microservice Spring Boot définit une configuration CORS restrictive
limitant les origines autorisées à l’application Angular~:

\begin{itemize}[leftmargin=*,nosep]
  \item \textbf{Origines autorisées}~: \texttt{http://localhost:4200} (configurable via \texttt{CORS\_ALLOWED\_ORIGIN})
  \item \textbf{Méthodes autorisées}~: \texttt{GET, POST, PUT, PATCH, DELETE, OPTIONS}
  \item \textbf{En-têtes autorisés}~: \texttt{*} (Authorization, Content-Type inclus)
  \item \textbf{Credentials}~: \texttt{allowCredentials=true} (requis pour les cookies de refresh)
\end{itemize}

\subsection{Désactivation de CSRF}\label{sec:csrf}

La protection CSRF est explicitement désactivée dans tous les microservices
(\texttt{csrf.disable()}). Cette décision est justifiée par le modèle
d’authentification sans cookie~: les requêtes API portent le JWT dans l'en-tête
\texttt{Authorization: Bearer}, et non dans un cookie de session. Un attaquant
ne peut pas exploiter le CSRF sans accès au token Bearer, qui n’est jamais envoyé
automatiquement par le navigateur.

\subsection{HikariCP et isolation de la base de données}

L'accès à Oracle 23ai est exclusivement réalisé via le pool de connexions HikariCP
configuré avec un compte dédié \texttt{gerai} (isolation du schéma). Les paramètres
de connexion sont externalisés via des variables d’environnement
(\texttt{DB\_URL}, \texttt{DB\_USER}, \texttt{DB\_PASSWORD}), jamais codés en dur.

\begin{table}[H]
\centering
\caption{Paramètres de sécurité HikariCP}
\label{tab:hikari_security}
\begin{tabularx}{\textwidth}{VlVlVLV}
\rowcolor{headerblue!55}
\hline
\textbf{Paramètre} & \textbf{Valeur} & \textbf{Justification sécurité} \\
\hline
\texttt{initialization-fail-timeout} & \texttt{-1} & Démarrage sans DB disponible (résilience) \\
\hline
\texttt{connection-timeout} & \texttt{5 000 ms} & Évite les blocages sur DB lente \\
\hline
\texttt{maximum-pool-size} & \texttt{10} & Limite la charge sur Oracle \\
\hline
\texttt{ddl-auto} & \texttt{none} & Protège le schéma en production \\
\hline
\end{tabularx}
\end{table}

% ─────────────────────────────────────────────────────────────
\section{Sécurisation du canal WebSocket STOMP}
% ─────────────────────────────────────────────────────────────

\subsection{Défi de sécurité WebSocket}

Le protocole HTTP standard ne s'applique pas à une connexion WebSocket après son
établissement~: une fois le handshake HTTP effectué, le canal devient TCP bidirectionnel.
Les en-têtes HTTP (y compris \texttt{Authorization}) ne sont donc plus disponibles pour
chaque message STOMP. Synapse résout ce problème via un intercepteur STOMP dédié.

\subsection{WebSocketAuthChannelInterceptor}

Le \texttt{WebSocketAuthChannelInterceptor} implémente
\texttt{ChannelInterceptor} et intercepte chaque frame STOMP de type
\texttt{CONNECT} pour en extraire et valider le JWT~:

\begin{enumerate}[leftmargin=*,nosep]
  \item Le client Angular envoie le JWT dans l'en-tête STOMP \texttt{Authorization}
        lors de la frame \texttt{CONNECT}.
  \item L'intercepteur extrait le token, supprime le préfixe \texttt{Bearer~} et
        appelle le \texttt{JwtDecoder} Spring pour valider la signature et l'expiration.
  \item Si valide, l’authentification Spring Security est propagée via
        \texttt{SecurityContextHolder}~; sinon, la connexion STOMP est rejetée.
\end{enumerate}

\subsection{Isolation des conversations}

Chaque conversation du chat est identifiée par un UUID stocké en base Oracle.
Les topics STOMP suivent le schéma \texttt{/topic/conversation/\{id\}}~: un
utilisateur qui tenterait de s'abonner à un topic dont il n’est pas participant
ne recevrait aucun message (la logique d'envoi côté serveur filtre les destinataires
par \texttt{participant\_id} avant toute publication).

% ─────────────────────────────────────────────────────────────
\section{Prévention des vulnérabilités OWASP}
% ─────────────────────────────────────────────────────────────

Le tableau~\ref{tab:owasp} ci-dessous détaille la couverture de Synapse vis-à-vis
des dix principales vulnérabilités OWASP Top 10 (édition 2021).

\begin{table}[H]
\centering
\caption{Couverture OWASP Top 10 dans Synapse}
\label{tab:owasp}
\begin{tabularx}{\textwidth}{VlVLVLV}
\rowcolor{headerblue!55}
\hline
\textbf{Vulnérabilité OWASP} & \textbf{Risque} & \textbf{Mesure implémentée dans Synapse} \\
\hline
A01 — Broken Access Control & Accès non autorisé aux données & RBAC via \texttt{@PreAuthorize} + JWT, guard Angular \\
\hline
A02 — Cryptographic Failures & Exposition de données sensibles & JWT signé RS256~; mots de passe jamais stockés (Keycloak) \\
\hline
A03 — Injection SQL & Destruction/vol de données & JPA/Hibernate (requêtes paramétrées)~; la requête native de l'analytics-service utilise des paramètres nommés \texttt{:param}, jamais de concaténation \\
\hline
A04 — Insecure Design & Architecture vulnérable & Session stateless~; PKCE~; Zero Trust inter-services \\
\hline
A05 — Security Misconfiguration & Exposition des configs & Variables d’environnement~; pas de credentials en dur \\
\hline
A06 — Vulnerable Components & Dépendances obsolètes & Spring Boot 3.5 (CVE patchés)~; Keycloak 26 \\
\hline
A07 — Auth Failures & Brute force, sessions volées & Keycloak gère les tentatives~; JWT court (5 min) + refresh \\
\hline
A08 — Data Integrity Failures & Tampering des données & JWT signé~; Oracle constraints CHECK \\
\hline
A09 — Security Logging & Absence de traçabilité & Table \texttt{AUDIT\_LOGS} Oracle~: chaque action (CREATE / UPDATE / DELETE / LOGIN / EXPORT) est journalisée avec l'identité Keycloak, l'adresse IP, le module et les valeurs avant/après \\
\hline
A10 — Server-Side Request Forgery & Proxy malveillant & Appels Groq côté serveur uniquement~; URLs fixes \\
\hline
\end{tabularx}
\end{table}

\subsection{Protection contre l'injection SQL}

Synapse utilise exclusivement Spring Data JPA avec des requêtes JPQL paramétrées
et des méthodes de repository dérivées. Les requêtes natives présentes dans
l'\textit{analytics-service} (pour la vue \texttt{V\_ATTRITION\_DATA}) sont
construites avec des paramètres nommés \texttt{:param}, jamais par concaténation de
chaînes. Les contraintes CHECK d'Oracle sur les valeurs de statut constituent une
deuxième ligne de défense contre les insertions de valeurs invalides.

\subsection{Protection contre le XSS}

Angular 21 protège nativement contre les attaques XSS par un mécanisme de
\textit{sanitization} automatique~: toutes les interpolations de templates
(\texttt{{{ }}}), les bindings \texttt{[innerHTML]} et les attributs sont
sanitisés avant rendu DOM via le \texttt{DomSanitizer}. Aucune donnée utilisateur
n’est injectée dans le DOM via \texttt{innerHTML} non sanitisé.

\subsection{Protection contre les attaques de type Broken Object Level Authorization}

Pour chaque endpoint sensible, Synapse vérifie que l’utilisateur est bien le
propriétaire de la ressource demandée~: le \texttt{sub} du JWT est résolu en
\texttt{employee\_id} Oracle via la méthode \texttt{findByKeycloakSub()}, et cet
identifiant est comparé à l'\texttt{employee\_id} de la ressource avant tout accès.
Un employé ne peut pas consulter les demandes d’un autre employé, même en devinant
l'identifiant.

% ─────────────────────────────────────────────────────────────
\section{Gestion des données sensibles}
% ─────────────────────────────────────────────────────────────

\subsection{Externalisation des secrets}

Tous les secrets (mots de passe, clés API, credentials SMTP) sont externalisés
via des variables d’environnement et des valeurs par défaut de développement.
Aucun secret n’est présent dans le code source versionné~:

\begin{table}[H]
\centering
\caption{Gestion des secrets par service}
\label{tab:secrets}
\begin{tabularx}{\textwidth}{VlVLVlV}
\rowcolor{headerblue!55}
\hline
\textbf{Secret} & \textbf{Variable d’environnement} & \textbf{Service(s) concerné(s)} \\
\hline
Mot de passe Oracle & \texttt{DB\_PASSWORD} & Tous les services \\
\hline
Mot de passe admin Keycloak & \texttt{KEYCLOAK\_ADMIN\_PASSWORD} & employe-service \\
\hline
Clé API Groq & \texttt{GROQ\_API\_KEY} & chat-service, projets-service \\
\hline
Credentials SMTP & \texttt{SMTP\_USER}, \texttt{SMTP\_PASS} & notification-service \\
\hline
Password Oracle master & \texttt{ORACLE\_PASSWORD} & Docker Compose \\
\hline
\end{tabularx}
\end{table}

\subsection{Données personnelles et RGPD}

Synapse stocke des données personnelles (nom, prénom, date de naissance, numéro CIN,
adresse, données médicales pour les congés longue maladie). Les mesures suivantes
réduisent le risque d'exposition~:

\begin{itemize}[leftmargin=*,nosep]
  \item \textbf{Accès segmenté par rôle}~: les données médicales (champs
        \texttt{MED\_APPROVED}, \texttt{MED\_COMMENT}) ne sont accessibles qu'aux
        gestionnaires et administrateurs RH (\texttt{RH}, \texttt{ADMIN\_RH}).
  \item \textbf{Pas de journalisation des données sensibles}~: les logs Spring Boot
        sont configurés sans affichage des valeurs de paramètres JPA (\texttt{show-sql=false}
        en production).
  \item \textbf{Chiffrement en transit}~: les API REST et les connexions WebSocket
        sont destinées à être servies via HTTPS en environnement de production.
  \item \textbf{Durée de conservation}~: les données de notifications et d'historique
        sont conservées sans purge automatique dans la version actuelle~; l'ajout
        d'un job Oracle de purge périodique (ex.~: notifications de plus de 6 mois)
        constitue une évolution prévue.
  \item \textbf{Droit à l'effacement~:} la désactivation d'un employé (\texttt{statut = INACTIF})
        révoque son accès Keycloak et masque ses données dans les interfaces~; la
        suppression physique des données personnelles (droit RGPD à l'effacement)
        n'est pas automatisée dans la version actuelle et constitue un axe d'amélioration.
\end{itemize}

\subsection{Gestion sécurisée des fichiers uploadés}

Les pièces jointes (justificatifs de congés, CV candidats, documents administratifs)
sont stockées sur le système de fichiers local (configurable via \texttt{FILE\_UPLOAD\_DIR})
et non dans Oracle. En environnement de production, \texttt{FILE\_UPLOAD\_DIR} doit
pointer vers un volume persistant (montage Docker ou stockage réseau) pour garantir
la survie des fichiers en cas de redémarrage du conteneur. Le \texttt{FileStorageService} valide~:
\begin{itemize}[leftmargin=*,nosep]
  \item La taille du fichier (limite~: 20 Mo par requête).
  \item Le type MIME via l'extension (les extensions non whitelistées sont rejetées).
  \item Lábsence de \textit{path traversal} dans le nom du fichier (nettoyage du
        nom original via \texttt{StringUtils.cleanPath()}).
\end{itemize}

\subsection{Traçabilité et journalisation des audits}

Synapse maintient une table \texttt{AUDIT\_LOGS} dans Oracle~23ai qui constitue
le registre central de toutes les actions effectuées sur la plateforme.
Chaque enregistrement capture l'identité Keycloak de l'auteur (\texttt{user\_id}),
l'action réalisée (\texttt{CREATE}, \texttt{UPDATE}, \texttt{DELETE}, \texttt{READ},
\texttt{LOGIN}, \texttt{LOGOUT}, \texttt{EXPORT}), le type et l'identifiant de
l'entité concernée, ainsi que l'état complet avant et après la modification
(\texttt{old\_values} et \texttt{new\_values} en JSON CLOB), l'adresse IP et le
\textit{user-agent} du client.

\begin{table}[H]
\centering
\caption{Structure de la table \texttt{AUDIT\_LOGS}}
\label{tab:audit_logs}
\begin{tabularx}{\textwidth}{VlVlVLV}
\rowcolor{headerblue!55}
\hline
\textbf{Colonne} & \textbf{Type} & \textbf{Utilité} \\
\hline
\texttt{user\_id} & UUID & Identifiant Keycloak de l'auteur de l'action \\
\hline
\texttt{action} & VARCHAR & Type d'opération~: \texttt{CREATE / UPDATE / DELETE / READ / LOGIN / EXPORT} \\
\hline
\texttt{entity\_type} & VARCHAR & Entité touchée~: \texttt{EMPLOYEE}, \texttt{DEMANDE}, \texttt{CONTRAT}, \texttt{PROJET}… \\
\hline
\texttt{old\_values} & JSON CLOB & État complet de l'entité avant modification \\
\hline
\texttt{new\_values} & JSON CLOB & État complet de l'entité après modification \\
\hline
\texttt{ip\_address} & VARCHAR & Adresse IP du client ayant émis la requête \\
\hline
\texttt{module} & VARCHAR & Microservice à l'origine de l'entrée \\
\hline
\texttt{created\_at} & TIMESTAMP & Date et heure précise de l'action \\
\hline
\end{tabularx}
\end{table}

Cette traçabilité répond à trois exigences complémentaires~: la conformité RGPD
(identifier qui a accédé ou exporté des données personnelles), la sécurité
forensique (reconstituer le fil des événements après un incident), et
l'auditabilité métier (prouver qu'une décision RH a bien été prise par la
bonne personne au bon moment). Les colonnes JSON \texttt{old\_values} et
\texttt{new\_values} permettent de rejouer ou d'inverser n'importe quelle
modification de manière programmatique.

% ─────────────────────────────────────────────────────────────
\section{Conclusion}
% ─────────────────────────────────────────────────────────────

Ce chapitre a présenté l’architecture de sécurité multicouche de Synapse~:
authentification déléguée à Keycloak, contrôle d’accès déclaratif à six rôles
cohérent entre frontend et backend, sessions stateless JWT, canal WebSocket
sécurisé et journalisation forensique via \texttt{AUDIT\_LOGS}. L’ensemble
garantit la traçabilité des actions et la conformité RGPD de la plateforme.
Le chapitre suivant présente la stratégie de tests et la validation des
fonctionnalités livrées par les sept sprints.
